\documentclass{article}

% if you need to pass options to natbib, use, e.g.:
%     \PassOptionsToPackage{numbers, compress}{natbib}
% before loading neurips_2026

% The authors should use one of these tracks.
% Before accepting by the NeurIPS conference, select one of the options below.
% 0. "default" for submission
 \usepackage{neurips_2026}
% the "default" option is equal to the "main" option, which is used for the Main Track with double-blind reviewing.
% 1. "main" option is used for the Main Track
%  \usepackage[main]{neurips_2026}
% 2. "position" option is used for the Position Paper Track
%  \usepackage[position]{neurips_2026}
% 3. "dandb" option is used for the Datasets & Benchmarks Track
 % \usepackage[dandb]{neurips_2026}
% 4. "creativeai" option is used for the Creative AI Track
%  \usepackage[creativeai]{neurips_2026}
% 5. "sglblindworkshop" option is used for the Workshop with single-blind reviewing
 % \usepackage[sglblindworkshop]{neurips_2026}
% 6. "dblblindworkshop" option is used for the Workshop with double-blind reviewing
%  \usepackage[dblblindworkshop]{neurips_2026}

% After being accepted, the authors should add "final" behind the track to compile a camera-ready version.
% 1. Main Track
 % \usepackage[main, final]{neurips_2026}
% 2. Position Paper Track
%  \usepackage[position, final]{neurips_2026}
% 3. Datasets & Benchmarks Track
 % \usepackage[dandb, final]{neurips_2026}
% 4. Creative AI Track
%  \usepackage[creativeai, final]{neurips_2026}
% 5. Workshop with single-blind reviewing
%  \usepackage[sglblindworkshop, final]{neurips_2026}
% 6. Workshop with double-blind reviewing
%  \usepackage[dblblindworkshop, final]{neurips_2026}
% Note. For the workshop paper template, both \title{} and \workshoptitle{} are required, with the former indicating the paper title shown in the title and the latter indicating the workshop title displayed in the footnote.
% For workshops (5., 6.), the authors should add the name of the workshop, "\workshoptitle" command is used to set the workshop title.
% \workshoptitle{WORKSHOP TITLE}

% "preprint" option is used for arXiv or other preprint submissions
 % \usepackage[preprint]{neurips_2026}

% to avoid loading the natbib package, add option nonatbib:
%    \usepackage[nonatbib]{neurips_2026}

\usepackage[utf8]{inputenc} % allow utf-8 input
\usepackage[T1]{fontenc}    % use 8-bit T1 fonts
\usepackage{hyperref}       % hyperlinks
\usepackage{url}            % simple URL typesetting
\usepackage{booktabs}       % professional-quality tables
\usepackage{amsfonts}       % blackboard math symbols
\usepackage{nicefrac}       % compact symbols for 1/2, etc.
\usepackage{microtype}      % microtypography
\usepackage{xcolor}         % colors

% Note. For the workshop paper template, both \title{} and \workshoptitle{} are required, with the former indicating the paper title shown in the title and the latter indicating the workshop title displayed in the footnote. 
\title{Sample-Efficient Learning of Control Using Kolmogorov–Arnold Networks}


% The \author macro works with any number of authors. There are two commands
% used to separate the names and addresses of multiple authors: \And and \AND.
%
% Using \And between authors leaves it to LaTeX to determine where to break the
% lines. Using \AND forces a line break at that point. So, if LaTeX puts 3 of 4
% authors names on the first line, and the last on the second line, try using
% \AND instead of \And before the third author name.


\author{%
  David S.~Hippocampus\thanks{Use footnote for providing further information
    about author (webpage, alternative address)---\emph{not} for acknowledging
    funding agencies.} \\
  Department of Computer Science\\
  Cranberry-Lemon University\\
  Pittsburgh, PA 15213 \\
  \texttt{hippo@cs.cranberry-lemon.edu} \\
  % examples of more authors
  % \And
  % Coauthor \\
  % Affiliation \\
  % Address \\
  % \texttt{email} \\
  % \AND
  % Coauthor \\
  % Affiliation \\
  % Address \\
  % \texttt{email} \\
  % \And
  % Coauthor \\
  % Affiliation \\
  % Address \\
  % \texttt{email} \\
  % \And
  % Coauthor \\
  % Affiliation \\
  % Address \\
  % \texttt{email} \\
}


\begin{document}


\maketitle


\begin{abstract}
  The abstract paragraph should be indented \nicefrac{1}{2}~inch (3~picas) on
  both the left- and right-hand margins. Use 10~point type, with a vertical
  spacing (leading) of 11~points.  The word \textbf{Abstract} must be centered,
  bold, and in point size 12. Two line spaces precede the abstract. The abstract
  must be limited to one paragraph.
\end{abstract}

\section{Introduction}

%%% PARAGRAPH 01: RL-based control (using Neural Networks) is great and outperforms
Reinforcement learning (RL) has become a powerful paradigm for learning feedback controllers directly from interaction with a system \cite{sutton2018reinforcement}.
Classical control methods, including model-based and model predictive control, can be highly effective when accurate dynamics models and tractable objectives are available, but these assumptions are often difficult to satisfy in complex, nonlinear, or noisy environments.
Value-based RL methods such as Q-learning instead, learn an action-value function that estimates the long-term return of taking an action in a given state \cite{watkins1992qlearning}.
Deep RL extends this idea by replacing tabular value functions or policies with neural networks, enabling controllers to approximate high-dimensional and nonlinear state-action relationships \cite{mnih2015humanlevel,schulman2017ppo}.
This flexibility has made neural RL attractive for control tasks in robotics, networking, and other cyber-physical systems where the controller must cope with sensor noise, nonlinear dynamics, delayed effects, and changing operating conditions.

%%% PARAGRAPH 02: What are challenges to RL-based control?
A central challenge in RL-based control is that learning a useful controller can require a large number of environment interactions.
Each sample in a control problem is not merely a labeled example, but a transition generated by running a policy in a simulator, a digital twin, or a physical system.
When these rollouts are computationally expensive, slow, unsafe, or subject to hardware wear, sample inefficiency becomes a practical barrier to deployment.
RL also requires balancing exploration and exploitation, designing reward functions that encode the intended behavior, and maintaining stability while the data distribution changes with the policy.
For safety-critical control, an additional limitation is that standard neural policies are difficult to inspect, verify, or translate into controller equations.
These issues motivate function approximators that can learn strong policies from fewer interactions while retaining enough structure to support analysis and interpretation.

%%% PARAGARAPH 03: TALK ABOUT KOLOMOGOROV ANROLD NETWORKS (DONT CHANGE THIS PARAGRAPH)
Kolmogorov-Arnold Networks (KANs) are a recently proposed class of neural network architectures inspired by the Kolmogorov-Arnold representation theorem, which states that any continuous multivariate function can be expressed as a finite superposition of univariate functions and summation operators \cite{kolmogorov1956representation, arnold1963representation}.
Unlike traditional neural networks such as multilayer perceptrons (MLPs), which rely on linear transformations followed by fixed nonlinear activation functions, KANs replace scalar weights with learnable univariate functions defined along the edges of the network \cite{liu2024kan}.
This formulation allows the model to directly learn functional relationships between variables rather than approximating them only through weighted sums of fixed activation functions.
As a result, KANs offer improved interpretability and significantly higher parameter efficiency while retaining strong approximation capabilities.
Recently, KANs have gained considerable attention, particularly in areas where understanding the underlying functional structure of the data is important \cite{somvanshi2025survey,essahraui2025kanoverview,sohail2024training}.
KANs are therefore particularly well suited for applications such as symbolic regression, scientific discovery, and physics-informed modeling, where both data efficiency and model transparency are essential \cite{somvanshi2025survey,essahraui2025kanoverview}.
These properties have also motivated the use of KANs in RL-based control, including PPO-based continuous control, robotic policy learning, load balancing, and explainable RL \cite{kich2024online,bayeh2026kanppo,chen2025kanpolicy,singh2026loadbalancing,dalto2025explainable}.

%%% PARAGRAPH 04: TALK ABOUT RESEARCH OF KANS REGARDING SAMPLE-EFFICIENCY --> TODO: make this shorter, and move the research gap at the end of this paragraph!
Existing evidence on the sample efficiency of KANs is promising but unsettled.
The original KAN paper reports that smaller KANs can match or outperform larger MLPs in several data-fitting and scientific-computing examples and suggests favorable neural scaling behavior \cite{liu2024kan}.
However, this evidence is closer to parameter and scaling efficiency than to a controlled sample-efficiency comparison.
In direct low-data comparisons, Pourkamali-Anaraki finds that MLPs with individualized learnable activations can outperform KANs when only a few hundred samples are available \cite{pourkamali2025lowdata}.
Similarly, Zeng et al. show that KANs do not uniformly dominate MLPs on irregular or noisy functions, even when parameter count and sample size are controlled \cite{zeng2024irregular}.
More positive evidence appears in structured settings: multifidelity KANs reduce the need for expensive high-fidelity data, KAN autoencoders perform well for low-data industrial fault detection, and QuadKAN reports faster learning and improved stability for quadruped locomotion \cite{howard2024multifidelity,villagomez2025kanae,wang2025quadkan}.
KAN-Dreamer provides a more cautious RL result, finding that FastKAN can match MLP sample efficiency and wall-clock performance in selected world-model components, while direct replacement of visual encoders or actor-critic networks can degrade efficiency \cite{shi2025kandreamer}.
Theoretical work also gives sample-complexity guarantees for residual KANs under smoothness and noiseless-sample assumptions, but these results do not settle empirical RL control settings with noisy rewards and nonstationary data \cite{kratsios2025besov}.

%%% PARAGRAPH 05: MOTIVATION FOR RESEARCH QUESTION
These mixed findings leave a clean research gap: it remains unclear when KANs improve sample efficiency over MLPs in RL-based control under comparable experimental budgets.
In this work, we study the sample-efficiency capabilities of KANs relative to MLPs in RL-based control.
We focus on the practical setting in which interaction data are expensive and compare architectures under controlled sample budgets, parameter budgets, and training protocols.
% Our study separates three notions that are often conflated in the KAN literature: parameter efficiency, wall-clock efficiency, and environment-sample efficiency.
% By doing so, we aim to clarify whether the structured univariate functions learned by KANs provide a useful inductive bias for control, or whether their additional optimization cost outweighs their benefits in common RL training regimes.
% This framing also supports a more careful use of KANs in control applications where data, compute, and interpretability constraints must be considered together.


% %%% PARAGRAPH 01: RL-based control (using Neural Networks) is great and outperforms
% Reinforcement Learning has demonstrated capabilties in learning to control complex systems, that are superior to conventional control approaches (model based control MPS, others?). 
% Effectively, what does it to? Learn the value function that maps the best (reward maximizing) response (action) to given (sensor) state of a system.
% Explain a bit, that first there was Q Learning, then neural network-based approaches to approximate the value function. and neural networks can caputre much more complex relationships, enabling them to work so much better than conventional control.
% RL-based NN control can caputre more non-linearity, more complexity in relationships, and also learn to cope with sensor noise.

% %%% PARAGRAPH 02: What are challenges to RL-based control?
% A central challenge in RL-based control is the learning aspect (sample inefficiency, balancing exploration exploitation, designing effective reward functions, and interpretability/transparency for formal verification and safety guarantees).
% (Why?) e.g. depending on the context
% My answer: because often computationally-expensive models and simulations are used to control complex systems for training, and that takes a lot of time and resources.

% %%% PARAGARAPH 03: TALK ABOUT KOLOMOGOROV ANROLD NETWORKS (DONT CHANGE THIS PARAGRAPH)
% Kolmogorov--Arnold Networks (KANs) are a recently proposed class of neural network architectures inspired by the Kolmogorov--Arnold representation theorem, which states that any continuous multivariate function can be expressed as a finite superposition of univariate functions and summation operators \cite{kolmogorov1956representation, arnold1963representation}. 
% Unlike traditional neural networks such as multilayer perceptrons (MLPs), which rely on linear transformations followed by fixed non-linear activation functions (e.g., ReLU or sigmoid), KANs replace scalar weights with learnable univariate functions defined along the edges of the network \cite{liu2024kan}. 
% This formulation allows the model to directly learn functional relationships between variables rather than approximating them through weighted sums of fixed activation functions. 
% As a result, KANs offer improved interpretability and significantly higher parameter efficiency while retaining strong approximation capabilities. 
% Recently, KANs have gained considerable attention, particularly in areas where understanding the underlying functional structure of the data is important \cite{somvanshi2025survey,essahraui2025kanoverview,sohail2024training}. 
% KANs are therefore particularly well suited for applications such as symbolic regression, scientific discovery, and physics-informed modelling, where both data efficiency and model transparency are essential \cite{somvanshi2025survey, essahraui2025kanoverview}. 
% This is, why KANs have been increasingly been explored in the context of RL-based control: 
% % Kolmogorov-Arnold Networks for Online Reinforcement Learning
% % https://ieeexplore.ieee.org/document/10773080/.
% % KAN Policy: Learning Efficient and Smooth Robotic Trajectories via Kolmogorov-Arnold Networks
% % https://ieeexplore.ieee.org/abstract/document/11151197
% % Interpretable Reinforcement Learning for Load Balancing using Kolmogorov-Arnold Networks
% % https://ieeexplore.ieee.org/abstract/document/11366542
% % KAN-Based On-Policy PPO Method for Continuous Reinforcement Learning
% % https://link.springer.com/chapter/10.1007/978-3-032-13589-6_1
% % Kolmogorov-Arnold Networks for Explainable Reinforcement Learning
% % https://www.diva-portal.org/smash/record.jsf?pid=diva2%3A2022559&dswid=6038

% %%% PARAGRAPH 04: TALK ABOUT RESEARCH OF KANS REGARDING SAMPLE-EFFICIENCY
% My prompt: is there any work related to KANs studying that KANs have higher sample-efficiency, meaning they can learn more from fewer sampels?

% ChatGPT answer: I’ll narrow this to “sample efficiency” rather than general performance: papers that explicitly vary training-set size, data budget, or environment interactions and compare KANs to MLPs.
% Yes — there is work directly related to KAN sample efficiency, but the evidence is mixed. The most careful answer is:
% KANs have been studied for low-data/sample-efficiency, but there is no settled result that KANs generally learn more from fewer samples than MLPs. Some KAN variants show better low-data behavior in structured domains, while fair comparisons often find MLPs match or outperform KANs.
% Most relevant papers
% Paper	Relevance to sample efficiency	Finding
% “Kolmogorov-Arnold Networks in Low-Data Regimes”	Most direct paper on “fewer samples”	This paper explicitly studies KANs vs MLPs in low-data regimes. It actually finds that MLPs with individualized learnable activations outperform KANs, especially around ~100 samples; in one additive-manufacturing task, MLP median accuracy was 0.91 vs KAN 0.53 with default hyperparameters.
% Original KAN paper: “KAN: Kolmogorov-Arnold Networks”	Early evidence for efficiency/scaling	Shows smaller KANs can match or beat larger MLPs in data fitting and PDE solving, and claims faster neural scaling laws. This is closer to parameter/scaling efficiency than clean sample-efficiency, but it motivates the idea.
% “KAN versus MLP on Irregular or Noisy Functions”	Controls both parameter count and training-sample size	Finds KAN does not always outperform MLP; increasing training samples helps, but both architectures struggle when noise obscures structure.
% “KAN-Dreamer”	Directly studies RL/world-model sample efficiency	Explicitly evaluates sample efficiency, wall-clock time, and asymptotic performance in DreamerV3. Result: FastKAN as reward/continue predictor gives parity with MLP sample efficiency and training speed, not a strong win.
% “QuadKAN”	RL locomotion; claims improved sample efficiency	Uses spline-parameterized KAN policy components for quadruped control and claims improved sample efficiency, smoother actions, and better returns vs baselines. This is domain-specific and not a broad MLP-vs-KAN theorem.
% “Comparative Evaluation of KAN Autoencoders… with Varying Training Set Sizes”	Varies training-set size in fault detection	Reports KAN autoencoders reaching strong fault-detection rates with fewer samples than the baseline autoencoder, e.g. EfficientKAN-AE reaching ≥90% FDR with 500 samples.
% “Multifidelity KANs”	Expensive-data setting	Uses low-fidelity models plus small high-fidelity datasets to reduce the amount of expensive high-fidelity data needed. This is highly relevant if your motivation is costly samples/experiments.
% “Kolmogorov-Arnold Networks: Approximation and Learning Guarantees…”	Theory/sample complexity	Gives approximation and sample-complexity guarantees for residual KANs under Besov-regularity/noiseless-sample assumptions. Useful theoretically, but not a broad empirical MLP comparison.
% My read
% The best paper to cite if you want to discuss low-data KANs is Pourkamali-Anaraki 2024, but it is a cautionary citation because it says KANs can degrade in low-data regimes compared with better-designed MLPs.
% The best positive evidence is more domain-specific: scientific computing/PDEs, multifidelity modeling, process monitoring, and some RL control tasks. The strongest defensible claim is:
% KANs and KAN variants can be more sample-efficient in structured problems where the target function has smooth, compositional, low-dimensional, or piecewise-smooth structure; however, systematic comparisons do not yet support the broad claim that KANs are generally more sample-efficient than MLPs.

% A clean research gap remains: KAN vs MLP under fixed compute, fixed parameter budget, fixed wall-clock, and varying sample budgets, especially in expensive-sampling settings like RL, active learning, Bayesian optimization, scientific simulations, and noisy reward/data regimes.

% MY THOUGH: Great, talk a bit about related works, then determine this as research gap.

% %%% PARAGRAPH 05: MOTIVATION FOR RESEARCH QUESTION
% In this work, we explore the sample-efficiency capabilities of KANs over MLPs in the context of reinforcement-learning based control. 

\section{Related Works}

\textbf{Kolmogorov--Arnold networks and learning guarantees.}
KANs were introduced as an alternative to MLPs in which learnable univariate functions are placed on edges rather than fixed nonlinearities on nodes \cite{liu2024kan}.
This architecture has motivated rapid follow-up work because it connects neural approximation to interpretable function decompositions and can use fewer parameters in some scientific-computing and symbolic-regression settings \cite{somvanshi2025survey,essahraui2025kanoverview}.
Theoretical work has begun to formalize these intuitions by proving approximation and sample-complexity guarantees for residual KANs on smooth function classes, including Besov spaces \cite{kratsios2025besov}.
These guarantees are important, but they are typically stated under assumptions that differ from RL control, where targets are noisy, bootstrapped, and induced by a changing policy.

\textbf{Low-data and fixed-sample empirical comparisons.}
Empirical evidence for KAN sample efficiency remains mixed.
The original KAN experiments show favorable scaling and strong performance with compact models, but they do not isolate environment-sample efficiency under fixed compute and fixed interaction budgets \cite{liu2024kan}.
Pourkamali-Anaraki directly studies low-data regimes and finds that MLPs with individualized learnable activations can outperform KANs with substantially fewer parameters, especially when only a few hundred samples are available \cite{pourkamali2025lowdata}.
Zeng et al. further show that KANs and MLPs trade advantages across regular, irregular, discontinuous, and noisy functions, with KANs converging quickly in some cases but failing to dominate across stabilized test loss \cite{zeng2024irregular}.
These comparisons suggest that KANs should not be treated as universally more sample-efficient than MLPs.

\textbf{Structured and expensive-data settings.}
More positive results appear when the problem structure matches the inductive bias of KAN edge functions.
Howard et al. propose multifidelity KANs that combine abundant low-fidelity data with limited high-fidelity data, reducing the amount of expensive high-fidelity data needed in scientific machine learning \cite{howard2024multifidelity}.
Villagomez and Mahalec evaluate KAN autoencoders for fault detection across varying training-set sizes and find that several KAN variants match or surpass an MLP-based orthogonal autoencoder with substantially less data \cite{villagomez2025kanae}.
These results are highly relevant for control domains with costly simulation or measurement, but they are supervised or unsupervised learning studies rather than online RL studies.

\textbf{KANs in reinforcement learning and control.}
Several recent works integrate KANs into RL pipelines.
Kich et al. replace MLP actor and critic components in PPO with KANs and report comparable performance with fewer parameters on DeepMind Control Proprio robotics benchmarks \cite{kich2024online}.
Bayeh et al. similarly propose KAN-PPO and MultKAN-PPO for continuous control in MuJoCo HalfCheetah, finding robust performance across several configurations while noting sensitivity to dropout design \cite{bayeh2026kanppo}.
D'Alto studies KAN-based PPO agents for CartPole and radio-access-network control, showing that KAN agents can match MLP performance with far fewer parameters and can be converted into sparse symbolic policies \cite{dalto2025explainable}.
Singh et al. apply a KAN actor in PPO for load balancing and emphasize the ability to extract interpretable controller equations \cite{singh2026loadbalancing}.
Together, these works support the feasibility of KAN-based RL, but they primarily emphasize parameter efficiency and interpretability rather than controlled sample-efficiency comparisons.

\textbf{KANs for robotic policies and world models.}
KANs have also been used in policy-learning systems outside the standard actor-critic replacement setting.
Chen et al. introduce KAN Policy for diffusion-based visuomotor policy learning and report smoother trajectories, shorter execution times, and improved success rates relative to diffusion-policy baselines \cite{chen2025kanpolicy}.
Wang and Tao propose QuadKAN for vision-guided quadruped locomotion and report faster convergence, improved returns, lower action jitter, and better robustness in structured gait-control settings \cite{wang2025quadkan}.
Shi and Luan study KAN-Dreamer by replacing selected DreamerV3 components with KAN or FastKAN modules, finding that FastKAN is promising for reward and continue prediction but that direct replacement of visual encoders or actor-critic modules can harm sample efficiency \cite{shi2025kandreamer}.
This component-level evidence reinforces the need to distinguish where KANs are useful in an RL stack rather than assuming that every MLP can be replaced beneficially.

\textbf{Position of this work.}
The literature therefore motivates a narrower and more testable question: whether KANs improve environment-sample efficiency for RL-based control when compared against MLPs under matched training protocols.
Our work focuses on this question by separating sample efficiency from parameter count and wall-clock speed.
This distinction is essential because a controller can use fewer parameters yet require more training time, or converge in fewer interactions while being slower per update.
By evaluating KANs and MLPs under controlled budgets, we aim to identify the conditions under which KANs are a useful inductive bias for control rather than only a compact or interpretable alternative.

\iffalse
%%%%%%%%%%%% THIS IS NOT RELEVANT ANYMORE
\section{Old Introduction}

Kolmogorov--Arnold Networks (KANs) are a recently proposed class of neural network architectures inspired by the Kolmogorov--Arnold representation theorem, which states that any continuous multivariate function can be expressed as a finite superposition of univariate functions and summation operators \cite{kolmogorov1956representation, arnold1963representation}. 
Unlike traditional neural networks such as multilayer perceptrons (MLPs), which rely on linear transformations followed by fixed non-linear activation functions (e.g., ReLU or sigmoid), KANs replace scalar weights with learnable univariate functions defined along the edges of the network \cite{liu2024kan}. 
This formulation allows the model to directly learn functional relationships between variables rather than approximating them through weighted sums of fixed activation functions. 
As a result, KANs offer improved interpretability and significantly higher parameter efficiency while retaining strong approximation capabilities. 
Recently, KANs have gained considerable attention, particularly in areas where understanding the underlying functional structure of the data is important \cite{somvanshi2025survey,essahraui2025kanoverview,sohail2024training}. 
KANs are therefore particularly well suited for applications such as symbolic regression, scientific discovery, and physics-informed modelling, where both data efficiency and model transparency are essential \cite{somvanshi2025survey, essahraui2025kanoverview}. 

The major challenges when working with KANs are speed and instability during training.
First, their slower training and inference, since univariate function evaluations are more expensive than matrix multiplications (as used in MLPs) and are not readily parallelisable on modern GPU architectures~\cite{liu2024kan,somvanshi2025survey,essahraui2025kanoverview}. 
Second, training is unstable; initialisation and hyper-parameters become critical, as poor starting values for spline coefficients can lead to vanishing gradients or erratic function behaviour~\cite{tan2019vanishing}.
These challenges currently limit the scalability of KANs and motivate the development of improved training strategies~\cite{sohail2024training}.

Recent works have proposed a multitude of approaches to these optimisation challenges, differing in the class of edge functions used, algorithm, loss function, and scope of optimization (from edge functions to network structure).
Despite this growing body of work, efficient and reliable training of KAN models remains an open research problem, particularly for high-dimensional data and large-scale machine learning tasks.





\newpage
\section{Old Related Works}
KANs provide competitive accuracies when compared to MLP architectures while requiring a fraction of parameters.
However, they are more sensitive to initialization and training schemes; KANs can benefit from adaptive optimizers and low learning rates~\cite{sohail2024training}.
Therefore, a growing branch of the literature resulted in a variety of approaches to train KANs, differing in the class of edge functions used, algorithm, loss function, and scope of optimization (from edge functions to network structure).

% %%%% OFF TOPIC
% ~\cite{sohail2024training} finds:
% Our study demonstrates that KANs provide competitive accuracies when compared to MLP architectures, but are more
% sensitive to choice of training scheme. For optimal performance KAN architectures must be trained using adaptive
% optimizers and low learning rates. Without such training schemes, KANs tend to overfit and exhibit high generalization
% error (compared to MLPs trained with the same scheme).
%%%%

\textbf{Gradient-Descent-Based Training.}
The Kolmogorov-Arnold representation theorem proposes a network structure with a depth of two layers, and does not impose any condition on the class of edge function.
\textit{Liu et al.}~\cite{liu2024kan} proposed KANs with more depth and suggest the use of \texttt{B-splines}~\cite{fakhoury2022exsplinet} as edge function, which enables training the degree parameters of the splines using gradient descent with back-propagation optimisers, such as adaptive moment estimation (ADAM), Broyden–Fletcher–Goldfarb–Shannon (BFGS), or stochastic gradient descent (SGD).
This however, leads to training instability and highly non-convex optimization problems, that are sensitive to knot placement.
% Gradient-Descent-Based Training with B-Splines ~\cite{liu2024kan}
%     Using B-Splines (Basis Splines) of degree k
%     Method: Back-propagation, using ADAM / LBFGS / SGD on spline coefficients
%     Issues: 
%         - highly non-convex optimization
%         - sensitive to know placement
%         - training instability
%\textbf{Basis Function Variation.}
Building on this, works suggested the use of alternative classes of edge functions, such as \texttt{wavelets}~\cite{bozorgasl2405wav}, \texttt{radial base functions} (RBF), \texttt{Chebyshev polynomials}~\cite{ss2024chebyshev}, \texttt{Gottlieb polynomials} or hybrid architectures with splines~\cite{ta2024bsrbf} to improve stability, enable faster convergence, and to effectively capture a multitude of frequency components.
However, these alternative basis functions introduce additional design choices and hyper-parameters, such as kernel widths, wavelet scales, or polynomial orders, which may increase model complexity and sensitivity to hyper-parameter tuning.
Moreover, some basis functions sacrifice the interpretability associated with spline-based representations and can incur higher computational cost due to more expensive function evaluations.
% Basis Function Variation
%     Wavelent-based KANs~\cite{bozorgasl2405wav}
%         Edge functions are expanded in wavelet bases.
%         Advantages:
%         multiresolution representation
%         faster convergence
%         better capture of high-frequency signals
%     RBF-based / Hybrid KANs~\cite{ta2024bsrbf}
%         Combine B-splines + radial basis functions.
%         Advantages
%         better smoothness
%         improved stability

\textbf{Physics-Informed Training.}
\textit{Rigas et al.}~\cite{rigas2024adaptive} suggest the integration of physics constraints in the form of partial differential equations (PDE) into the loss function, which can significantly increase the training speed in the context of physics-informed neural networks (PINNs and PIKANs). 
This approach however, is limited to physics-contextual applications only.
% Physics-Informed Training~\cite{rigas2024adaptive}
%     KANs are being trained using PDE constraints, similar to PINNs.
%     Instead of only data loss:
%     L=Ldata+λLphysics
%     Contributions:
%         adaptive grid refinement
%         faster training (up to 84× faster than original KAN code)
%     Applications:
%         PDE solving
%         scientific computing

\textbf{Multilevel Training.}
To address the optimisation difficulties arising from spline parameterisation, \textit{Southworth et al.}~\cite{southworth2026multilevel} propose a multilevel training strategy inspired by multigrid methods in numerical analysis.
The key idea is to train the network in a coarse-to-fine manner by initially using a low-resolution spline grid and progressively refining the spline-knot placement during training.
At each refinement step, parameters from the coarser model are interpolated to initialise the finer representation, which stabilises optimisation and improves convergence speed.
Such hierarchical refinement reduces the effective dimensionality of the optimisation problem during early training stages and mitigates poor local minima caused by high-resolution spline parameterisations.
% Multilevel Training ~\cite{southworth2026multilevel}
%     Train the network coarse-to-fine by gradually refining the spline grid.
%     Main contribution:
%         hierarchical training scheme using spline knot refinement
%         analytic interpolation between model levels


\textbf{Hierarchical and Non-Backpropagation Training.}
An alternative line of research seeks to avoid gradient-based optimisation entirely.
\textit{Dudek et al.}~\cite{dudek2025hkan} introduce the Hierarchical Kolmogorov–Arnold Network (HKAN), which trains the model sequentially without back-propagation.
In this approach, basis functions are fixed and layers are trained one at a time by solving convex least-squares problems.
This decomposition simplifies optimisation and improves numerical stability, although it may limit the flexibility of the learned representation due to the absence of end-to-end optimisation.
% Hierarchical Training~\cite{dudek2025hkan}
%     Recent work explores training without backpropagation.
%     Example: Hierarchical KAN (HKAN) 
%     Idea:
%         Fix basis functions randomly
%         Train layers sequentially
%         Solve convex least-squares problems
%     Advantages
%         no backprop
%         more stable training
%         faster computation


\textbf{Training via Nonlinear Equation Solvers.}
Another perspective formulates KAN training as a nonlinear system identification problem.
\textit{Poluektov et al.}~\cite{poluektov2025construction} propose solving for the network parameters using a Newton–Kaczmarz iterative scheme.
This method directly estimates the parameters of the edge functions by iteratively solving systems of non-linear equations derived from the network representation.
Such approaches can achieve fast local convergence but typically require careful initialisation and may scale poorly with network size.
% Nonlinear Equation Methods / Newton-Kaczmarz method~\cite{poluektov2025construction}
%     Another line treats training as nonlinear system solving.
%     Method
%     Newton–Kaczmarz iterative solver.
%     Advantages
%         locally fast convergence
%         direct parameter estimation

\textbf{Architecture-Adaptive Training.}
Finally, several works explore training strategies that adapt the network structure itself.
For instance, Progressive KAN (proKAN)~\cite{gyanchandani2024prokan} dynamically grows the architecture during training by incrementally adding functional blocks.
Model complexity is increased only when required, which helps mitigate overfitting and reduces the risk of training overly expressive networks.
These approaches highlight that, beyond the optimisation of edge functions, the structure of KANs itself can play an important role in training stability and generalisation.
% Architecture-Training
%     Progressive KAN (proKAN)~\cite{gyanchandani2024prokan}
%         Idea:
%         grow architecture dynamically
%         stop adding blocks when overfitting appears
        
\newpage
\fi

\section{Methods}

>>>>Ideas for making training more efficient:
(1) Adaptive Spline-Knot-Position Learning (and not only Spline-Coefficients)
(2) Sparse Edge Learning (Edge Pruning) --> sparsity penalty to loss function, so that more neurons turn to zero
(3) Curriculum Training --> start with simple edge functions, train them, and then enhance complexity slowly
(4) Low-Rank KAN --> factorise functions that share structure, to reduce function evaluations

>>>>A completely different training paradigm:
We train with MLP (stable, fast) and then convert to KAN.
Currently approach do "distillation": So in practice, if you want “MLP → KAN”, you currently use distillation: generate a dataset (or use the training data plus extra synthetic queries), feed it through the trained MLP, and train a KAN to regress to those outputs. This is function projection in the learning sense, not an analytic algebraic conversion.

{\color{cyan}KC: Latter sounds great to me, if you need a recipe I guess this one works and is open-sourced and only needs few hours of training on 8x H200: \url{https://arxiv.org/pdf/2601.22156}. Something very basic could be just comparing strong-to-weak distillation from scratch for a KAN-based and MLP-based model, to see if the KAN can be more accurate, but I guess this has been done already?}

\section{Experiments}

\section{Conclusion}

\newpage


%%%%%%%%%%%%%%%%%%%%%%%%%%%%%%%%%%%%%%%%%%%%%%%%%%%%%%%%%%%%
\newpage
% \section*{References}

% References follow the acknowledgments in the camera-ready paper. Use unnumbered first-level heading for
% the references. Any choice of citation style is acceptable as long as you are
% consistent. It is permissible to reduce the font size to \verb+small+ (9 point)
% when listing the references.
% Note that the Reference section does not count towards the page limit.
% \medskip

\bibliographystyle{unsrt}
\bibliography{references}


%%%%%%%%%%%%%%%%%%%%%%%%%%%%%%%%%%%%%%%%%%%%%%%%%%%%%%%%%%%%

\appendix

\section{Technical appendices and supplementary material}
Technical appendices with additional results, figures, graphs, and proofs may be submitted with the paper submission before the full submission deadline (see above). You can upload a ZIP file for videos or code, but do not upload a separate PDF file for the appendix. There is no page limit for the technical appendices. 

Note: Think of the appendix as ``optional reading'' for reviewers. The paper must be able to stand alone without the appendix; for example, adding critical experiments that support the main claims to an appendix is inappropriate. 

%%%%%%%%%%%%%%%%%%%%%%%%%%%%%%%%%%%%%%%%%%%%%%%%%%%%%%%%%%%%

% \newpage
% \input{checklist.tex}


\end{document}
